\scanpage{108}

\uline{Genus of an entire function}

It happens that one has $\sum\frac{1}{|a_n|^{k+1}}<+\infty$ for some $k$, (let $k$ be the smallest integer satisfying it) and moreover that in the repre-\newline sentation (2), $g(z)$ is a polynomial -- incidentally determined up to a constant.

Let $r$ be its degree, the \deleted{smaller} \inserted{larger} of the numbers $r,k$ is called the genus of $f(z)$.
\inserted{It is an \uline{integer} $\geq 0$.}

Thus if $f(z)$ is of genus $\leq r$, then one has the formula
\[
f(z)=e^{\alpha_0+\alpha_1 z\mathord{\cdot}\mathord{-}+\alpha_n z^n}\prod_n\left(1-\frac{z}{a_n}\right)e^{\frac{z}{a_n}+\mathord{-}\mathord{\cdot}\frac{z^{\unsure{n|r}}}{r a_n}}\tag*{(3)}
\]

To say that $f(z)$ is of genus zero, means that $f(z)$ is of the form
\[
\begin{aligned}\vphantom{f(z)}\\[-2pt]\text{with}\quad\end{aligned}
\left\{\begin{aligned}
f(z)&=\mathrm c^{\mathrm{te}}\times\prod_n\left(1-\frac{z}{a_n}\right)\unsure{.|,}\\[-2pt]
\sum_n\frac{1}{|a_n|}&<+\infty
\end{aligned}\right.\tag*{(4)}
\]

\uline{Order of an entire function} \quad Put $M(r)=\mathop{\rm Sup}_{z\unsure{=|\leq}r}|f(z)|$

Then $\rho=\overline{\lim}\frac{\log\log M(r)}{\substack{\deleted{\unsure{z}}\\\inserted{\log r}}}$ \quad $(0\leq\rho\leq+\infty)$, $\rho$ is called the order of $f(z)$, and $f(z)$ is said to be of order
\par\noindent\hspace{2cm}\unsure{$=$}\par\noindent
finite if $\rho<+\infty$. This means that there exists $k>0$ such that $|f(z)|\leq e^{\unsure{(z\vert^k)|(\lvert z\rvert^k)}}$. Then $\rho$ is the \deleted{\unsure{smallest number|smallest value}} \inserted{infimum of the} $\sigma>0$ such that

\margin{109}
